% Numerical experiment plan for the NASA Space Apps 2026 entry "Sea Level Detective"
% Same visual style as the argument tree. Compile with pdflatex (twice).
\documentclass[10pt]{article}
\usepackage{mathptmx}\usepackage[scaled=0.9]{helvet}\usepackage[T1]{fontenc}
\usepackage[margin=0.5in]{geometry}
\usepackage{amsmath,amssymb}
\usepackage[most]{tcolorbox}
\usepackage{booktabs,tabularx,enumitem,xcolor,longtable}
\usepackage{xurl}
\usepackage[hidelinks]{hyperref}
\usepackage{microtype}
\setlist{nosep,leftmargin=1.2em}
\newcolumntype{Y}{>{\raggedright\arraybackslash}X}
\setlength{\parindent}{0pt}

\tcbset{evbase/.style={enhanced,breakable,sharp corners,boxrule=0.4pt,colframe=black!70,colback=white,
  fonttitle=\sffamily\bfseries\footnotesize,fontupper=\small,left=4pt,right=4pt,top=2pt,bottom=2pt,
  boxsep=1pt,toptitle=1pt,bottomtitle=1pt,before skip=4pt,after skip=4pt}}
\newtcolorbox{nubox}[1][NUMERICAL EXPERIMENT]{evbase,colbacktitle=black!20,coltitle=black,title={#1}}
\newtcolorbox{cabox}[1][CAVEATS AND BOUNDARIES]{evbase,colback=black!4,colbacktitle=black!4,coltitle=black,
  borderline={0.4pt}{0pt}{black!70,dashed},boxrule=0pt,title={#1}}
\newtcolorbox{expbox}[1]{enhanced,sharp corners,boxrule=0pt,colback=white,colframe=white,
  borderline north={1.2pt}{0pt}{black},borderline south={0.4pt}{0pt}{black},
  colbacktitle=white,coltitle=black,title={#1},fonttitle=\sffamily\bfseries\large,
  fontupper=\small\itshape,left=0pt,right=0pt,top=1pt,bottom=4pt,toptitle=4pt,bottomtitle=0pt,
  before skip=8pt,after skip=4pt}
% one row of an experiment card
\newcommand{\row}[2]{\textbf{\sffamily\footnotesize #1} & #2\\}
\newcommand{\card}[1]{\begin{nubox}[]\begin{tabularx}{\linewidth}{@{}p{2.3cm} Y@{}}#1\end{tabularx}\end{nubox}}

\begin{document}
\begin{center}
{\sffamily\LARGE\bfseries Numerical Experiment Plan: Sea Level Detective}\\[4pt]
{\small NASA Space Apps 2026 \textperiodcentered{} built around the argument tree (claims C1--C5) \textperiodcentered{} status 2026-11-14}\\[2pt]
{\footnotesize Rule: every experiment is placed next to a published number (the \emph{anchor}). An anchor mismatch requires diagnosis
(different data, period, corrections or method), not automatic classification as a bug or a finding.}
\end{center}
\vspace{-4pt}\noindent\rule{\linewidth}{0.8pt}

\begin{tcolorbox}[enhanced,sharp corners,boxrule=0.6pt,colframe=black,colback=black!5,colbacktitle=black,coltitle=white,
  title={OVERVIEW},fonttitle=\sffamily\bfseries\footnotesize,fontupper=\small,left=5pt,right=5pt,toptitle=1pt,bottomtitle=1pt]
\begin{tabularx}{\linewidth}{@{}l l Y l l@{}}
\toprule
ID & Claim & Question & Priority & Time\\\midrule
E0 & C1.1 & Does our NASA-SSH global series match NASA's own GMSL indicator, and what rate and acceleration result? & must & 3 h\\
E1 & C3.1--3.2 & Do our tests keep a 5\% false-positive rate, and how many years do they need? (synthetic, no login) & must & 2 h\\
E2 & C2.1--2.2 & Does the El Ni\~no regression show the west--east seesaw? (two boxes + map) & must & 2 h\\
E3 & C2.3, C3.3 & How much do 10-year trends scatter, raw vs ENSO-removed? (descriptive) & must & 1 h\\
E4 & C4.1--4.2 & ENSO at U.S. tide gauges (no login); box vs gauge agreement (with NASA data) & must & 3 h\\
E5 & C2.4 & Does ENSO removal reproduce the 2010--11 dip and the vanishing ``slowdown''? & should & 1 h\\
E6 & C1.2 & What share of the ocean has a significant local trend once natural variability counts? & should & 2 h\\
E7 & C5 & Do 1997--98, 2015--16 and 2023--24 Kelvin waves travel at theoretical speed? & could & 2 h\\
E8 & C4.4 & How large was the 2023 Panama rainfall deficit in GPM? & could & 1 h\\\bottomrule
\end{tabularx}
\textbf{Status 2026-11-14 19:20 UTC:} S (pipeline, data standard, 29 unit tests passing), E1 and the gauge part of E4 are \textbf{done}
(\texttt{pipeline/results/summary.md}). E0, E2, E3 and the E4 satellite part wait for an Earthdata token.
Must is about 11 h of analysis, plus setup and download (S, about 3 h including retries). Run the parts that need no login (E1, E4 gauges)
first, while the NASA files download. The pipeline is in \texttt{pipeline/}; \texttt{run\_all.py --no-nasa} runs those parts.
\end{tcolorbox}

\begin{expbox}{D. NASA datasets available (checked in NASA's CMR catalogue, 2026-11-14)}
Coverage as listed in CMR. All PO.DAAC and GES DISC data need a free Earthdata login. \emph{Use} names the experiments that need each dataset.
\end{expbox}
\begin{nubox}[]
\footnotesize
\begin{tabularx}{\linewidth}{@{}l l l Y l@{}}
\toprule
Tier & Short name (provider) & Coverage & What it is & Use\\\midrule
core & \texttt{NASA\_SSH\_REF\_SIMPLE\_GRID\_V1} (PO.DAAC) & 1992-10 -- now & SSH anomaly, 0.5\textdegree, every 7 days, reference missions & all\\
core & \texttt{NASA\_SSH\_GMSL\_INDICATOR} (PO.DAAC) & 1993-01 -- now & NASA global mean sea level from the same grid & E0, E5\\
core & \texttt{NASA\_SSH\_ENSO\_INDICATOR} (PO.DAAC) & 1993-01 -- now & NASA ENSO index from sea-surface height & E1--E6\\
core & \texttt{NASA\_SSH\_PDO\_INDICATOR} (PO.DAAC) & 1993-01 -- now & NASA Pacific Decadal Oscillation index from SSH & E3\\
support & \texttt{MERGED\_TP\_J1\_OSTM\_OST\_GMSL\_ASCII\_V52} (PO.DAAC) & 1992-09 -- & MEaSUREs multi-mission GMSL (independent processing) & E0 check\\
support & \texttt{TELLUS\_GRAC-GRFO\_MASCON\_CRI\_GRID\_RL06.3\_V4} (PO.DAAC) & 2002-04 -- now & GRACE/GRACE-FO mass (ocean + land water) & E5 mechanism\\
support & \texttt{HOMAGE\_STERIC\_OHC\_TIME\_SERIES\_v01} (PO.DAAC) & 1978 -- & global steric height / ocean heat content & E5 mechanism\\
optional & \texttt{GPM\_3IMERGM} v07 (GES DISC) & 1998-01 -- 2025-09 & monthly precipitation, 0.1\textdegree & E8\\
optional & \texttt{CCMP\_WINDS\_10M6HR\_L4\_V3.1} (PO.DAAC) & 1993-01 -- now & 6-hourly 10 m winds, 0.25\textdegree & E7 (forcing)\\
optional & \texttt{SWOT\_L2\_LR\_SSH\_BASIC\_D} (PO.DAAC) & 2022-12 -- now & 2 km swath SSH & site visual only\\
optional & \texttt{M2TMNXSLV} (GES DISC) & 1980-01 -- now & MERRA-2 monthly surface fields (pressure, wind) & E4 (East Coast)\\
not used & OSCAR, MUR SST, SMAP SSS (PO.DAAC) & various & currents, SST, salinity & out of scope\\\bottomrule
\end{tabularx}
\end{nubox}
\begin{cabox}[DATASET NOTES (FROM THE NASA-SSH INDICATORS USER GUIDE, READ 2026-11-14)]
\textbf{Indicators.} ENSO, PDO and IOD indicators are the second (``indicator'') mode of regional cyclostationary EOFs of the NASA-SSH grid,
computed over 1993--2019 after removing each cell's linear trend; the first mode is the seasonal cycle. The modes are projected onto
along-track SSH to give a time series to the present. They are ASCII files. Because they are built from the same SSH, regressing SSH on them
partly explains SSH by itself. We therefore report every ENSO-removed result twice: with the NASA ENSO indicator (the site default) and with
NOAA's SST-based ONI (independent). \textbf{GMSL indicator}: area-weighted mean of the same grid, in \textbf{cm}, 7-day steps, \textbf{not}
GIA-corrected (DOI 10.5067/NSIND-GMSV1).
\textbf{Access.} Every NASA file above needs a free Earthdata login (tested: PO.DAAC and GES DISC HTTPS and OPeNDAP return 401 or a login
redirect without one). No public mirror holds the long NASA-SSH record: NOAA ERDDAP's altimetry starts in 2017, and NASA's sea-level web pages
host no data files. NOAA tide gauges and ONI need no login.
\textbf{Other.} GPM IMERG Final monthly ends 2025-09 in CMR, which is enough for 2023. The GRACE CRI mascon gives the ocean-mass vs
land-water split that Boening et al.\ used for the 2010--11 dip.
\end{cabox}

\begin{expbox}{Shared setup (S)}
One pipeline, built once, used by every experiment.
\end{expbox}
\card{
\row{Data}{Core NASA set (table D below): the NASA-SSH grid (about 1{,}700 weekly files, about 2 GB), plus NASA's own GMSL, ENSO and PDO
sea-surface-height indicators, so the El Ni\~no removal uses a NASA index. NOAA ONI is a cross-check only. NOAA CO-OPS monthly mean sea level for the
Battery 8518750, La Jolla 9410230 and San Francisco 9414290 (public API) is the only non-NASA input.}
\row{Access}{Earthdata login is needed for every NASA file. Set \texttt{EARTHDATA\_TOKEN} in the environment (never in code), then run
\texttt{fetch\_nasa.py}. Alternatively, the professor runs the download with his token and shares \texttt{box\_monthly.csv} and the 1\textdegree{}
grid. Smoke-test one grid file first (variable \texttt{ssha} in metres; \texttt{latitude}, \texttt{longitude}).}
\row{Processing}{Weekly to monthly means (removes the 10-day window overlap). Area-weighted ($\cos\varphi$) box means.
Monthly climatology removed (deseasonalized). Output: one compact NetCDF/CSV of monthly series per box plus a 1\textdegree{} monthly grid.}
\row{Boxes}{Global ($\pm66$\textdegree), W.\ Pacific (5\textdegree S--5\textdegree N, 120\textdegree E--160\textdegree E), E.\ Pacific
(5\textdegree S--5\textdegree N, 150\textdegree W--90\textdegree W), California shelf, NY Bight shelf, Panama (rain only). Fixed
before looking at results.}
\row{Statistics}{Sen slope with CI; Hamed--Rao modified Mann--Kendall; seasonal Kendall (Hirsch et al.\ 1982); OLS with AR(1)-adjusted
CI as cross-check. ENSO removal: a \emph{joint} fit of trend + annual/semiannual harmonics + index at a lag fixed in advance (0; lags 0--6
shown only as a sensitivity table). Subtract only the fitted index component. Run with the NASA ENSO indicator and with ONI. Monthly
averaging does not remove autocorrelation, so the AR(1)-aware tests are always used.}
\row{Tools}{Python: xarray, netCDF4, earthaccess, numpy, scipy, statsmodels, pymannkendall. Laptop is enough (RTX 4060 not needed).}
}

% ---------------------------------------------------------------- E0
\begin{expbox}{E0. Reproduce the global rate and acceleration (C1.1) \hfill must}
If this fails, nothing downstream can be trusted.
\end{expbox}
\card{
\row{Method}{Global box, monthly, 1993-01 to 2017-12 (Nerem's window). Quadratic fit: rate at the window centre and $2\times$ the quadratic
coefficient. Repeat with ONI regressed out, and on 1993--latest.}
\row{Anchor}{Nerem et al.\ 2018, \emph{context only}: about 3 $\pm$ 0.4 mm/yr. Their acceleration is 0.097 mm/yr$^2$ before and 0.084 after
the Pinatubo and ENSO adjustments (multivariate EOF, not ONI regression). They apply GIA; the NASA GMSL indicator does not.}
\row{Check}{Validation against matched \texttt{NASA\_SSH\_GMSL\_INDICATOR} (same grid, cm converted to mm, matched timestamps): bias, RMSE and
trend difference, not correlation alone. Then report our rate and acceleration with CIs, uncorrected for GIA, next to Nerem's. Not a pass gate.}
\row{Site output}{Global curve with the fitted quadratic; the ``Global'' demo example.}
}

% ---------------------------------------------------------------- E1
\begin{expbox}{E1. Validate the trend tests on synthetic series (C3.1, C3.2) \hfill must}
Makes ``significant'' on the site mean what it says.
\end{expbox}
\card{
\row{Method}{Monte Carlo, $\geq$1{,}000 series per setting: trend $\omega$ + seasonal cycle + monthly AR(1) noise with \emph{marginal} SD
$\sigma_N$ and $\varphi$ (estimated from deseasonalized gauge residuals now, from the NASA boxes later). Lengths 10--35 years in 5-year steps;
$\omega \in \{0,\ldots,4\}$ mm/yr. Tests: naive OLS, MK, Hamed--Rao MK, seasonal Kendall, AR(1)-OLS. False-positive rates with binomial
95\% CIs; power curves; the 90\% power crossing interpolated. A separate run adds an ENSO-like term (empirical power only).}
\row{Anchor}{Weatherhead et al.\ 1998, \emph{plain AR(1) case only}: $n^* \approx \big[(3.3\,\sigma_N/|\omega|)\sqrt{(1+\varphi)/(1-\varphi)}\big]^{2/3}$
years (90\% power, 5\% level; monthly stationary AR(1), marginal SD). Hamed and Rao 1998: plain MK over-rejects under positive autocorrelation.}
\row{Read-out}{Which tests hold about 5\% at $\omega=0$ (naive OLS and plain MK are expected not to). AR(1)-OLS crossing next to $n^*$.}
\row{Site output}{``How many years do you need?'' panel per region; justifies the test the site uses.}
}

% ---------------------------------------------------------------- E2
\begin{expbox}{E2. The El Ni\~no seesaw (C2.1, C2.2) \hfill must}
The picture that makes the story.
\end{expbox}
\card{
\row{Method}{Per 1\textdegree{} cell, regress deseasonalized SSHA on ONI (fixed lag): map of mm per \textdegree C of ONI. W vs E Pacific
box series 1993--latest with El Ni\~no years shaded. Peak 2015/16 anomalies per box.}
\row{Anchor}{Alory and Delcroix 2002: zonal seesaw with its fulcrum at the dateline. Widlansky et al.\ 2017: western Pacific drops exceeding
30 cm in 2015/16 (island tide gauges).}
\row{Read-out}{Where the regression coefficient changes sign along the equator; the W and E box anomalies in 1997/98, 2015/16 and 2023/24,
next to the published island values (box means are expected to be smaller than island extremes). Descriptive, no gate.}
\row{Site output}{Seesaw map + the two-box time series: demo example 1, ``same cause, opposite trends''.}
}

% ---------------------------------------------------------------- E3
\begin{expbox}{E3. The short-window trap (C2.3, C3.3) \hfill must}
Quantifies ``short records mislead''.
\end{expbox}
\card{
\row{Method}{For every box: Sen slope over every 10-year window (sliding monthly) vs the full-record slope, raw and ENSO-removed. Overlapping
windows are dependent: the spread is descriptive, not a sample of independent replications.}
\row{Anchor}{Zhang and Church 2012, \emph{context only}: Jan 1993--Jun 2009, 50\textdegree S--60\textdegree N, 100\textdegree E--70\textdegree W,
high-pass MEI and low-pass PDO (25/37-month filters), 5-month smoothing; about 20\% vs almost 60\% of basin variance. Our box $R^2$ is not
comparable unless that setup is rebuilt (stretch goal).}
\row{Read-out}{Window-slope distributions per box, raw vs ENSO-removed, and their $R^2$ with and without the index. No gate.}
\row{Site output}{The window slider and the ``remove El Ni\~no'' toggle: demo example 4, ``the window trap''.}
}

% ---------------------------------------------------------------- E4
\begin{expbox}{E4. U.S. coasts: satellite vs tide gauge, ENSO amplitude (C4.1, C4.2) \hfill must}
Grounds the local panels; checks that 0.5\textdegree{} boxes mean something at the coast.
\end{expbox}
\card{
\row{Method}{Deseasonalized monthly series: California box vs La Jolla and San Francisco; NY Bight box vs the Battery. Correlation, trend
difference. ENSO regression amplitude at each gauge, peak 1997/98 and 2015/16 anomalies. Battery: cool-season (Oct--Apr) mean anomaly,
El Ni\~no winters vs others (Welch and permutation tests). Surge-day counts need daily data: deferred.}
\row{Anchor}{Hamlington et al.\ 2015: ENSO contributes up to 20 cm along the U.S. West Coast. Sweet and Zervas 2011: more cool-season surge
days in El Ni\~no at four East Coast gauges. Colle et al.\ 2010: no statistically established ENSO link for New York.}
\row{Read-out}{Peak ENSO-related anomalies at La Jolla/SF next to Hamlington's ``up to 20 cm'' (an upper envelope, not a target). Box--gauge
correlation as found; trend differences discussed as land motion (gauges) vs geocentric (altimetry). Battery comparison reported as found.}
\row{Site output}{California and New York demo panels with gauge overlay.}
}

% ---------------------------------------------------------------- E5--E8
\begin{expbox}{E5. Global dip and ``slowdown'' (C2.4) \hfill should}
Shows El Ni\~no matters even for the global number.
\end{expbox}
\card{
\row{Method}{Global series 2010--2012 drop (our grid vs \texttt{NASA\_SSH\_GMSL\_INDICATOR}); GRACE ocean mass over the same months for the
mechanism; trend 2005--2014 raw vs ENSO-removed.}
\row{Anchor}{Boening et al.\ 2012: about 5 mm drop, early 2010 to mid 2011. Cazenave et al.\ 2014: slowdown disappears after correction
(read the paper before using this anchor).}
\row{Read-out}{Size and timing of the 2010--11 drop next to Boening's 5 mm; GRACE mass change over the same months; 2005--2014 slopes, raw and corrected. No gate.}
}

\begin{expbox}{E6. Local significance map (C1.2) \hfill should}
How much of the ocean shows a real rise once natural variability is in the noise.
\end{expbox}
\card{
\row{Method}{Per 1\textdegree{} cell 1993--2019: Sen slope, Hamed--Rao $p$, with and without ENSO removal. Area fraction significant at 5\%.}
\row{Anchor}{Prandi et al.\ 2021: 98\% of the ocean rising significantly, but with measurement error only.}
\row{Read-out}{A separate descriptive map. Prandi uses 90\% confidence, a measurement-error covariance and other data and estimators, so the two
fractions are not directly comparable. Field significance is not claimed.}
}

\begin{expbox}{E7. Kelvin-wave speeds, three El Ni\~nos (C5) \hfill could}
Extends the professor's 2015--16 analysis; feeds the explainer page.
\end{expbox}
\card{
\row{Method}{Professor's slant-stack and $k$--$\omega$ code on 2\textdegree S--2\textdegree N for 1997--98, 2015--16, 2023--24 (weekly data).}
\row{Anchor}{Boulanger and Menkes 1995: Kelvin and first Rossby waves at theoretical speeds. Professor's 2015--16: +1.50 / +1.62 m/s.}
\row{Pass}{Eastward speeds consistent across events and with first-mode theory, or the difference explained (e.g.\ higher modes).}
}

\begin{expbox}{E8. Panama 2023 rainfall (C4.4) \hfill could}
Only if E0--E4 are done.
\end{expbox}
\card{
\row{Method}{GPM IMERG monthly over the canal watershed, 2001--latest; 2023 May--Dec anomaly vs climatology; ONI correlation of rainy-season totals.}
\row{Anchor}{World Weather Attribution 2024: low 2023 rainy-season rainfall; about 8\% less rain in an El Ni\~no year.}
\row{Pass}{2023 among the driest rainy seasons in the record; sign of the ONI correlation negative.}
}

\begin{cabox}[RULES FOR ALL EXPERIMENTS]
Boxes, lag and tests are fixed before results are seen. Every reported trend carries its window, $n$ (months), method and CI. Anchors are
context: a mismatch is diagnosed (data, period, corrections, method), never tuned away. ``Not significant'' is shown as such. No PINN retraining.
Plan revised after a feasibility review by Codex, 2026-11-14.
\end{cabox}

\begin{cabox}[SCHEDULE (SAT NIGHT -- SUN)]
Tonight: S (data + pipeline), E0, E1 in parallel (E1 needs no data). Sunday morning: E2, E3, E4. Midday: E5, E6 if on time; E7, E8 only
if the site is already working. 14:00 Sunday: freeze numbers, fill the argument tree's numerical boxes, update the narrative verdicts.
\end{cabox}
\end{document}
